\documentclass[../../../include/open-logic-section]{subfiles}

\begin{document}

\olfileid{sth}{spine}{zf}
\olsection{$\Z$, $\ZF$: ఒక మైలురాయి}

పునాదితో మరో ముఖ్యమైన మైలురాయికి చేరుకున్నాం. $\Zminus$, $\ZFminus$ సిద్ధాంతాలను పరిశీలించాం; వాటిలో ఏదో ఒకటి ``తీసివేయబడిందని'' చెప్పాం. తీసివేయబడినది పునాది. కనుక:

\begin{defn}
$\Z$ సిద్ధాంతం $\Zminus$కు పునాదిని చేర్చుతుంది. అందువల్ల దాని స్వయంసిద్ధ సూత్రాలు విస్తరణీయత, సమ్మేళనం, జతలు, ఘాత సమితులు, అనంతత్వం, పునాది, వేరుచేయు పథకపు అన్ని సందర్భాలు.

$\ZF$ సిద్ధాంతం $\ZFminus$కు పునాదిని చేర్చుతుంది. మరో మాటలో, $\ZF$ ప్రతిస్థాపన పథకపు అన్ని సందర్భాలను~$\Z$కు చేర్చుతుంది.
\end{defn}

అయినా మీకు ఒక ప్రశ్న రావచ్చు. $\ZFminus$లో నియమితత్వం పునాదికి సమానార్థకమై, నియమితత్వాన్ని సమర్థించడం స్పష్టమైతే, దాన్ని ప్రాథమిక స్వయంసిద్ధ సూత్రంగా తీసుకోకుండా పునాదిని ఎందుకు తీసుకోవాలి?

చారిత్రక కారణాలను (ఎవరు ఏది ఎప్పుడు రూపొందించారనే విషయాలను) పక్కనపెడితే, ప్రధాన కారణం ఇది: $V_\alpha$ల నిర్వచనం వాడకుండానే పునాదిని ప్రకటించవచ్చు. ఆ నిర్వచనం \olref[recursion]{sec}లోని మొత్తం కృషిపై ఆధారపడింది: అతిపరిమిత పునరావృత్తి సమర్థితమని చూపడానికి దాన్ని నిరూపించాల్సి వచ్చింది. కానీ ఆ నిరూపణలో \emph{ప్రతిస్థాపన} వాడాం. కనుక పునాది, నియమితత్వం $\ZFminus$ పరంగా సమానార్థకాలైనా, $\Zminus$ పరంగా అలా కావు.

వాస్తవానికి ఈ సాధారణ వ్యాఖ్య సూచించేదానికంటే విషయం మరింత తీవ్రమైనది. ఈ పుస్తక పరిధికి మించిన విషయమే అయినా, $\Zminus$, $\Z$ రెండూ $V_\alpha$లను నిర్వచించడానికి తగినంత బలమైనవి కావు. కాబట్టి మీరు $\Z$లో మాత్రమే పనిచేస్తుంటే, మనం రూపొందించిన నియమితత్వ వాక్యానికి \emph{అర్థమే ఉండదు}. అందుకే మన అధికారిక స్వయంసిద్ధ సూత్రం నియమితత్వం కాక పునాది.

ఇకపై, వేరుగా చెప్పనంతవరకు, మరే వ్యాఖ్య లేకుండానే $\ZF$లో పనిచేస్తాం.

\end{document}
